\chapter{Concurrence et avenir}\label{ch:fm:competition-and-the-future}

Lorsque le régulateur américain des marchés financiers (la SEC) proposa, en décembre 2022, de faire entrer en concurrence, dans des enchères, de nombreux ordres de détail, son propre communiqué constatait que les broker-dealers
envoyaient plus de 90~\% des ordres exécutables des investisseurs particuliers sur des actions cotées à six dealers hors bourse, les wholesalers, et que deux d'entre eux
traitaient environ 66~\% du volume des wholesalers. La proposition a été retirée en juin 2025. Le trading ressemble à un marché de nombreuses entreprises~; métier par métier, c'est souvent un marché de quelques-unes.
Ce dernier chapitre demande pourquoi, ce qui protège les quelques acteurs, et ce que font à la réponse trois changements en cours (actifs tokenisés,
la \omterm{def:fm:competition-and-the-future:rtc}{négociation en continu} et les outils d'apprentissage automatique).

\section{Avantages concurrentiels durables dans le trading}

\begin{definition}[Avantage concurrentiel durable]\label{def:fm:competition-and-the-future:moat}
Un \emph{avantage concurrentiel durable}\index{avantage concurrentiel durable} (economic moat) est un avantage durable qui permet à une entreprise de dégager des rendements supérieurs à son coût du capital dans une activité sans
que des concurrents entrent tant que ces rendements ne sont pas épuisés~: un avantage de coût lié à l'échelle, à une information propriétaire ou à un flux d'ordres, les \omterm{def:fm:build-against-buy:lockin}{coûts de changement} supportés par les clients,
des effets de réseau, ou une licence ou une relation que d'autres ne peuvent obtenir rapidement.
\end{definition}

Le livre~11 a montré qu'un edge de trading s'érode à mesure que d'autres le découvrent~: la part captée diminue quand des concurrents entrent. Les avantages durables ne sont donc que rarement les
signaux. Ce sont les coûts fixes qu'un nouvel entrant doit payer avant sa première transaction (la \omterm{def:fm:the-economics-of-a-trading-firm:fixed}{base de coûts fixes} du chapitre~1, la \omterm{def:fm:the-proprietary-market-making-firm:treadmill}{course technologique} du chapitre~2,
la \omterm{def:fm:technology-strategy:race}{course aux armements de la latence} du chapitre~19), le flux d'ordres qui passe par des relations et des contrats (chapitre~23), et les données et l'historique qu'une entreprise a
accumulés. Les travaux de Sutton sur les coûts irrécupérables et la structure de marché distinguent deux cas. Lorsque le coût fixe d'entrée est exogène, un marché plus grand
accueille davantage d'entreprises, et la concentration baisse à mesure qu'il grandit. Lorsque les entreprises peuvent dépenser pour être meilleures (plus rapides, avec de meilleures données, de meilleurs modèles) et que les clients
récompensent les meilleures, les dépenses elles-mêmes croissent avec le marché, et la concentration ne baisse pas forcément du tout.

\begin{definition}[Concentration du marché, indice de Herfindahl--Hirschman]\label{def:fm:competition-and-the-future:hhi}
La \emph{concentration du marché}\index{concentration du marché} est le degré auquel le volume ou le chiffre d'affaires d'une activité est détenu par quelques entreprises, mesuré par
des ratios de concentration (la part cumulée des $k$ plus grandes entreprises, $\mathrm{CR}_k$) ou par l'\emph{indice de Herfindahl--Hirschman}\index{indice de Herfindahl--Hirschman},
la somme des carrés des parts de marché en pourcentage,
\[ \mathrm{HHI}=\sum_{i=1}^n(100\,s_i)^2, \]
de valeurs proches de 0 pour de nombreuses petites entreprises à 10\,000 pour une seule (l'indice de Herfindahl des plateformes du livre~1, sur une échelle de 0 à 10\,000).
\end{definition}

\section{Consolidation et structure de marché}

Les sources publiques donnent rarement la part de chaque entreprise~; le communiqué en fournit deux. Ils encadrent l'indice.

\begin{proposition}[Bornes à partir de parts partielles]\label{prop:fm:competition-and-the-future:bounds}
Si une activité compte $n$ entreprises dont les $k$ plus grandes détiennent une part cumulée $S_k$, l'indice est minimal lorsque les $k$ leaders sont égaux et que les autres entreprises
sont égales entre elles, $\mathrm{HHI}_{\min}=10^4\bigl(S_k^2/k+(1-S_k)^2/(n-k)\bigr)$, pourvu que $(1-S_k)/(n-k)\le S_k/k$. Pour $n=6$, $k=2$ et $S_2=0.66$, elle
se situe entre 2\,467 et 3\,668.
\end{proposition}

\begin{proof}
À total fixé, une somme de carrés est minimale lorsque les termes sont égaux~; l'appliquer séparément aux leaders et aux autres donne le
minimum, qui est réalisable lorsque chaque autre entreprise n'est pas plus grande que chaque leader. Pour le maximum, on fixe la part $b$ du plus petit leader~: le plus grand
leader détient $S_2-b$, et les autres, chacun au plus $b$, détiennent $1-S_2$~; avec $b=(1-S_2)/4=0.085$ les quatre autres doivent tous être égaux à $b$, et le leader
détient 0.575, ce qui donne $10^4(0.575^2+5\times0.085^2)=3\,668$. Une recherche sur $b$ (\texttt{firm.moats.hhi\_bounds}) confirme qu'un $b$ plus grand donne un indice plus faible.
\end{proof}

\begin{omfigure}
\begin{tikzpicture}
\begin{axis}[width=10.5cm,height=5cm,ybar,bar width=10pt,ymin=0,ymax=65,ylabel={\small part (\%)},xtick={1,2,3,4,5,6},
  xlabel={\small wholesaler, par rang},tick label style={font=\footnotesize},grid=major,grid style={gray!20},table/col sep=comma,
  legend style={font=\scriptsize,at={(0.97,0.95)},anchor=north east},legend cell align=left]
\addplot[fill=omDef!60,draw=omDef] table[x=rank,y=low]{figdata/desk/30-competition-and-the-future/bounds.csv};
\addplot[fill=omThm!45,draw=omThm] table[x=rank,y=high]{figdata/desk/30-competition-and-the-future/bounds.csv};
\legend{le moins concentré (indice 2\,467),le plus concentré (indice 3\,668)}
\end{axis}
\end{tikzpicture}\par
\omcaption{Les deux vecteurs de parts qui bornent l'indice pour six wholesalers dont les deux plus grands détiennent 66~\%, soit les chiffres du communiqué de la SEC de 2022. Dans l'un ou l'autre
cas, l'activité est très concentrée au regard du seuil de 1\,800 des lignes directrices de 2023 sur les fusions. Données~: \texttt{fm\_competition.public\_bounds}.}\label{fig:fm:competition-and-the-future:bounds}
\end{omfigure}

\begin{dated}{2026-09}{Les sources publiques du chapitre}\label{dat:fm:competition-and-the-future:record}
\textbf{Wholesaling}~: la proposition de la SEC relative à la Order Competition Rule (Release 34-96495, publiée le 3 janvier 2023) constatait que plus de 90~\% des ordres exécutables des
investisseurs particuliers en actions NMS acheminés vers six wholesalers, dont deux concentraient environ 66~\% du volume d'actions exécutées par les wholesalers au premier
trimestre 2022~; la Commission a retiré la proposition le 17 juin 2025. \textbf{Lignes directrices sur les fusions} (DOJ et FTC, décembre 2023)~: un indice supérieur à 1\,800 est
très concentré~; une hausse supérieure à 100 dans un tel marché, ou une part fusionnée supérieure à 30~\% avec une hausse supérieure à 100, est présumée restreindre
sensiblement la concurrence. \textbf{\omterm{def:fm:competition-and-the-future:token}{Fonds tokenisé}}~: le prospectus d'un fonds monétaire gouvernemental américain enregistré, daté du 1\up{er} août 2026, indique que son
agent de transfert tient le registre officiel de propriété des actions dans un système intégré à une blockchain sur des réseaux publics. \textbf{\omterm{def:fm:competition-and-the-future:rtc}{Négociation en continu}}~: la SEC
a agréé la 24X National Exchange le 27 novembre 2024 pour négocier des actions cotées 23 heures par jour, cinq jours par semaine~; la NSCC a étendu ses horaires de compensation à
24x5 le 29 juin 2026, les bourses devant suivre d'ici la fin de 2026.
\end{dated}

\begin{definition}[Équilibre à entrée libre]\label{def:fm:competition-and-the-future:freeentry}
Un \emph{équilibre à entrée libre}\index{equilibre a entree libre@équilibre à entrée libre} (free-entry equilibrium) est le nombre d'entreprises auquel chaque entreprise d'une activité couvre son coût fixe, tandis qu'un
entrant supplémentaire ne le ferait pas~: $\pi(n)\ge F>\pi(n+1)$, où $\pi(n)$ est le résultat d'exploitation de chaque entreprise lorsque $n$ entreprises se font concurrence et $F$ est le coût fixe
d'exercice de l'activité.
\end{definition}

Sur un marché de Cournot à demande linéaire $a-bQ$, un coût marginal commun $c$ et $n$ entreprises identiques, chacune gagne $\pi(n)=S/(n+1)^2$ avec
$S=(a-c)^2/b$, donc le nombre d'entreprises à entrée libre vaut $n^*=\lfloor\sqrt{S/F}\rfloor-1$ et l'indice vaut $10\,000/n^*$. La concentration est fixée par le rapport entre la taille du
marché et le coût fixe, et non par l'un ou l'autre pris isolément.

\section{Tutoriel~: le coût fixe baisse}\label{tut:fm:competition-and-the-future:tut}

\begin{tutorial}
\textbf{Objectif.} Calculer le nombre d'entreprises à entrée libre et l'indice aux paramètres du livre~; diviser par deux puis doubler le coût fixe~; fusionner deux entreprises, avec
et sans synergies~; et encadrer l'indice à partir des chiffres publics. \textbf{État final~:}
les \cref{fig:fm:competition-and-the-future:entry,fig:fm:competition-and-the-future:merger}.
\begin{tutsteps}
\item \textbf{Paramètres.} $S=\$3\,000$ millions par an (ordonnée à l'origine de la demande 10 et coût marginal 4, en points de base), et un coût fixe $F=\$50$ millions
par an et par entreprise~; ces deux valeurs sont à titre illustratif.
\item \textbf{Entrée libre.} \texttt{free\_entry(3000, 50)} donne six entreprises, chacune gagnant 61.2 millions de dollars avant et 11.2 millions de dollars après son coût fixe~;
l'indice vaut 1\,667.
\item \textbf{Diviser par deux et doubler.} À 25 millions de dollars, neuf entreprises (1\,111)~; à 100 millions de dollars, quatre (2\,500).
\item \textbf{Le planificateur.} \texttt{planner} maximise le surplus des consommateurs plus les profits moins les coûts fixes~: trois entreprises à 50 millions de dollars.
\item \textbf{Une fusion.} Deux des six entreprises fusionnent et économisent un coût fixe~; trouver ensuite la synergie de coût marginal qui maintient le prix inchangé.
\end{tutsteps}
\end{tutorial}

\omcode{code/firm/moats/firm_moats.py}{33}{64}{Bornes de l'indice à partir de parts partielles, nombre d'entreprises à entrée libre, surplus total et
nombre retenu par le planificateur.}\label{lst:fm:competition-and-the-future:entry}

\begin{omfigure}
\begin{tikzpicture}
\begin{axis}[width=10.5cm,height=5.6cm,xlabel={\small coût fixe par entreprise (millions de dollars par an)},ylabel={\small indice (HHI)},xmin=10,xmax=200,ymin=0,
  ymax=6000,tick label style={font=\footnotesize},grid=major,grid style={gray!20},table/col sep=comma,legend cell align=left,
  legend style={font=\scriptsize,at={(0.03,0.97)},anchor=north west},scaled y ticks=false,
  yticklabel style={/pgf/number format/.cd,1000 sep={,}}]
\fill[omThm!12] (axis cs:10,2467) rectangle (axis cs:200,3668);
\addplot[omDef,very thick,const plot] table[x=F,y=hhi]{figdata/desk/30-competition-and-the-future/entry.csv};
\draw[gray,dashed] (axis cs:10,1800) -- (axis cs:200,1800) node[pos=0.8,below,font=\scriptsize] {1\,800~: très concentré};
\node[font=\scriptsize,omThm,anchor=west] at (axis cs:14,3070) {wholesalers, 2022~: bornes};
\draw[black] (axis cs:50,1666.7) circle[radius=2.5pt];
\legend{indice à entrée libre}
\end{axis}
\end{tikzpicture}\par
\omcaption{L'indice à entrée libre en fonction du coût fixe pour $S=\$3\,000$ millions~: six entreprises et 1\,667 à 50 millions de dollars (cerclé), neuf et 1\,111
à 25 millions de dollars, quatre et 2\,500 à 100 millions de dollars. La bande ombrée est la plage que permettent les chiffres publics des wholesalers. Données~:
\texttt{fm\_competition.entry}.}\label{fig:fm:competition-and-the-future:entry}
\end{omfigure}

Le modèle symétrique aux paramètres du livre donne six entreprises et un indice de 1\,667, sous le seuil des lignes directrices~; les chiffres publics des wholesalers
impliquent au moins 2\,467. L'écart tient à la symétrie du modèle~: les entreprises réelles diffèrent par leurs coûts, et une entreprise aux coûts plus faibles prend une part plus grande. Diviser par deux le
coût fixe, comme le font le cloud computing et les outils open source pour la recherche et l'infrastructure, fait passer le nombre à entrée libre de six à neuf~; le doubler,
comme le fait un nouveau palier de vitesse pour les entreprises qui doivent l'acheter, le ramène à quatre.

\begin{remark}[Entrée excédentaire]\label{rem:fm:competition-and-the-future:excess}
Le planificateur choisirait trois entreprises à 50~millions de dollars, et non six~: chaque entrant prend de l'activité aux autres, et son gain privé dépasse le
surplus qu'il ajoute. Le surplus total vaut 1\,169~millions de dollars en entrée libre, contre 1\,256~millions de dollars à trois entreprises. C'est le
résultat d'entrée excédentaire de Mankiw et Whinston dans le modèle du chapitre~; ce n'est pas un argument pour moins d'entreprises dans toute activité réelle, où les entrants apportent aussi de nouveaux produits et des coûts plus bas.
\end{remark}

La consolidation dans le trading a suivi l'arithmétique des coûts fixes. Knight Capital et GETCO ont fusionné au sein de KCG Holdings le 1\up{er} juillet 2013~; Virtu
Financial a racheté KCG le 20 juillet 2017 pour 20.00~\$ par action en numéraire, puis Investment Technology Group le 1\up{er} mars 2019 pour environ 1.0~milliard de dollars
(voir la \cref{fig:fm:competition-and-the-future:chain}). Chaque opération a réuni des entreprises dont les coûts fixes (technologie, connectivité, données, conformité)
se recoupaient~; le chapitre~20 a lu la seconde comme une décision de \omterm{def:fm:build-against-buy:bab}{faire ou acheter}.

\begin{omfigure}
\begin{tikzpicture}[font=\small,firm/.style={draw=omDef,fill=omDef!8,rounded corners,align=center,minimum height=0.8cm,text width=2.6cm},
  arr/.style={-{Stealth[length=2mm]},thick,omDef}]
\node[firm] (kn) at (0,1.4) {Knight Capital};
\node[firm] (ge) at (0,0) {GETCO};
\node[firm] (kcg) at (3.6,0.7) {KCG Holdings\\{\scriptsize 1\up{er} juillet 2013}};
\node[firm] (vi) at (9.2,0.7) {Virtu Financial\\{\scriptsize KCG~: 20 juillet 2017}};
\node[firm] (itg) at (9.2,-1.7) {Investment\\Technology Group};
\draw[arr] (kn) -- (kcg);
\draw[arr] (ge) -- (kcg);
\draw[arr] (kcg) -- node[above,font=\scriptsize] {20.00~\$ par action} (vi);
\draw[arr] (itg) -- node[right,font=\scriptsize,align=left] {1\up{er} mars 2019\\environ 1.0~milliard de dollars} (vi);
\end{tikzpicture}\par
\omcaption{Une chaîne de consolidation d'après les dépôts publics~: deux entreprises ont fusionné en 2013, l'entreprise issue de la fusion a été rachetée en 2017, et le repreneur a acquis
Investment Technology Group en 2019.}\label{fig:fm:competition-and-the-future:chain}
\end{omfigure}

\begin{omfigure}
\begin{tikzpicture}
\begin{axis}[width=10.5cm,height=5.2cm,xlabel={\small synergie de coût marginal de l'entreprise fusionnée (pb)},ylabel={\small prix (pb)},xmin=0,xmax=1.2,
  ymin=4.78,ymax=5.02,tick label style={font=\footnotesize},grid=major,grid style={gray!20},table/col sep=comma,
  yticklabel style={/pgf/number format/.cd,fixed,precision=2}]
\addplot[omDef,very thick] table[x=synergy,y=price]{figdata/desk/30-competition-and-the-future/merger.csv};
\draw[omThm,dashed] (axis cs:0,4.857) -- (axis cs:1.2,4.857) node[pos=0.05,below right,font=\scriptsize] {prix avant la fusion};
\draw[gray,dotted] (axis cs:0.857,4.78) -- (axis cs:0.857,5.02) node[pos=0.93,right,font=\scriptsize] {0.857};
\end{axis}
\end{tikzpicture}\par
\omcaption{Deux des six entreprises de Cournot fusionnent~: le prix après la fusion en fonction de la réduction du coût marginal de l'entreprise fusionnée. Sans elle, le prix
passe de 4.857 à 5.000 (2.9~\%)~; une réduction de 0.857, soit 21.4~\% du coût marginal, rétablit l'ancien prix. Données~:
\texttt{fm\_competition.merger}.}\label{fig:fm:competition-and-the-future:merger}
\end{omfigure}

\omcode{code/firm/moats/firm_moats.py}{67}{90}{Équilibre de Cournot à coûts inégaux, une fusion avec une synergie de coût marginal, et la synergie qui maintient
le prix inchangé.}\label{lst:fm:competition-and-the-future:merge}

La fusion dans le modèle montre pourquoi la consolidation a lieu même lorsqu'elle fait monter les prix. Sans synergie de coût marginal, le prix monte de 2.9~\% et
l'indice passe de 1\,667 à 2\,000, mais l'entreprise fusionnée économise un coût fixe~: le profit des deux entreprises après coûts fixes passe de 22.4~millions de dollars à 33.3
millions de dollars, tandis que les consommateurs perdent 60~millions de dollars de surplus. Avec la synergie qui maintient le prix, la part de l'entreprise fusionnée atteint 33.3~\% et l'indice 2\,222,
au-dessus des deux présomptions des lignes directrices, bien que les acheteurs ne paient pas davantage~: la concentration mesure la structure, non le préjudice, et c'est pourquoi les lignes directrices
traitent leurs seuils comme des présomptions à vérifier.

\section{Actifs tokenisés}

\begin{definition}[Tokenisation d'actifs, fonds tokenisé]\label{def:fm:competition-and-the-future:token}
La \emph{tokenisation d'actifs}\index{tokenisation d'actifs} enregistre la propriété d'un actif financier (une part de fonds, une obligation, un dépôt) sous forme de jetons sur une
blockchain, l'émetteur, son agent de transfert ou un \omterm{def:fm:the-asset-manager-and-the-fund:admin}{dépositaire} restant juridiquement responsable de l'enregistrement. Un
\emph{fonds tokenisé}\index{fonds tokenisé} est un fonds dont le registre officiel des parts est tenu, en tout ou partie, dans un tel système.
\end{definition}

Le fonds monétaire enregistré de l'encadré daté montre l'état actuel de l'enregistrement~: son agent de transfert tient le registre officiel dans un
système permissionné sur des blockchains publiques, avec une liste blanche de portefeuilles agréés dans un smart contract (livre~3), le pouvoir de corriger les erreurs et de limiter
les transferts, ainsi que les transferts pair à pair de parts entre portefeuilles agréés. Pour une société de trading, le changement est d'abord opérationnel avant d'être stratégique. Une
part pouvant circuler entre portefeuilles à toute heure peut servir de collatéral en dehors des heures bancaires, ce qui intéresse les marges et le financement du chapitre~14~; ses
registres sont tenus par un agent de transfert qui peut annuler des transferts, si bien que ce n'est pas un actif au porteur. Savoir si la tokenisation abaisse le coût fixe d'entrée dans une
activité (en supprimant les intermédiaires) ou l'augmente (nouvelle conservation, contrôles et connexions) est la question que le modèle ci-dessus rend précise.

\section{Marchés 24~h sur 24}

\begin{definition}[Négociation en continu (24~h sur 24)]\label{def:fm:competition-and-the-future:rtc}
La \emph{négociation en continu}\index{négociation en continu} étend les séances de négociation d'un marché à la plupart ou à la totalité des heures de la semaine, dans le cas des actions cotées,
à 23 ou 24~heures les jours ouvrés, la compensation, les données de marché et la surveillance nécessaires aux séances fonctionnant sur les mêmes plages.
\end{definition}

L'ordonnance de la Commission enregistrant 24X National Exchange décrit la structure~: les séances habituelles de pré-ouverture, de cœur de séance et de post-clôture, de 4~h à
19~h, et une séance de nuit de 20~h à 4~h les nuits précédant les jours ouvrés, les week-ends et les jours fériés ayant été retirés de la
proposition. La chambre de compensation a ouvert la voie~: la compensation de la NSCC court de dimanche 20~h à vendredi 20~h. Pour une entreprise, des horaires plus longs font monter les coûts fixes
(effectifs, supervision, réponse aux incidents la nuit, conception des astreintes du chapitre~21) alors que le volume de nuit est, au départ, faible. Dans le
modèle du chapitre, cela correspond à une hausse de $F$ pour toute entreprise qui choisit d'être présente, et à un $S$ plus petit par séance~: l'activité de nuit
soutiendra moins d'entreprises que le jour.

\section{Ce que changent les outils d'apprentissage automatique}

Les plateformes d'apprentissage automatique du livre~12 et les grands modèles de langage (LLM) réduisent deux coûts à la fois. Le coût de l'infrastructure de recherche baisse~: bibliothèques open source,
accélérateurs loués et modèles préentraînés font ce qu'une équipe construisait autrefois. Il en va de même pour le coût du travail courant (rapprochements, documentation, revue de code,
premières ébauches de recherche), l'organisation du chapitre~21 avec moins de personnes par fonction. Ces deux effets abaissent $F$, et le modèle à entrée libre prédit davantage d'entreprises. Le deuxième cas de Sutton
invite à aller plus loin~: si les outils permettent aussi aux meilleures entreprises de dépenser pour mieux faire (plus de données, modèles plus grands, itérations plus rapides), et si
l'activité récompense le fait d'être le meilleur, les dépenses augmentent avec le marché et la part des leaders n'a pas de raison de baisser. Les deux effets se distinguent dans
les chiffres propres à une entreprise~: quels coûts ont baissé (la \omterm{def:fm:the-economics-of-a-trading-firm:fixed}{base de coûts fixes} du chapitre~1), et si sa capture de marché a tenu face au meilleur concurrent ou face au concurrent moyen.

\begin{method}[Lire la structure d'une activité]\label{met:fm:competition-and-the-future:read}
\begin{enumerate}
\item Mesurer la concentration à partir de ce qui est public~: nombre d'acteurs, parts des premiers, bornes de l'indice.
\item Estimer la taille de l'activité et le coût fixe pour y être présent~; comparer le nombre d'entreprises à entrée libre avec celles qui sont observées.
\item Se demander quels coûts sont exogènes et lesquels les entreprises choisissent de dépenser pour être meilleures~; le second type entretient la concentration à mesure que le marché croît.
\item Pour chaque changement (un outil, une structure de marché, une règle), se demander s'il déplace $S$, $F$ ou la course aux dépenses, et pour qui.
\end{enumerate}
\end{method}

\section{Construction~: le modèle des avantages concurrentiels durables}\label{bld:fm:competition-and-the-future:moats}

\begin{build}
\bfield{Rôle} Concentration à partir de parts ou de parts partielles, nombre d'entreprises à entrée libre, nombre optimal du planificateur, et fusions de Cournot avec synergies de coût fixe-
et de coût marginal.
\bfield{Interface} \texttt{firm.moats}~: \texttt{hhi}, \texttt{cr}, \texttt{hhi\_bounds}, \texttt{free\_entry}, \texttt{welfare}, \texttt{planner},
\texttt{cournot}, \texttt{merge}, \texttt{synergy\_for\_price}, \texttt{scenarios}.
\bfield{Règles} Les parts sont normalisées~; une entreprise de Cournot sans production lève une erreur plutôt que de se retirer en silence~; les bornes rejettent les données de parts
qui ne peuvent pas être ordonnées.
\bfield{Tests d'acceptation} \texttt{code/firm/moats/tests/}~: l'indice sur des exemples faits à la main~; des bornes qui encadrent des vecteurs de parts aléatoires admissibles~; la
condition d'entrée libre et entrée excédentaire~; la hausse de prix après fusion et la synergie qui la supprime.
\bfield{Pour aller plus loin} Coûts irrécupérables endogènes (les entreprises choisissent la qualité à un coût)~; entrée avec des coûts asymétriques~; deux séances (jour et nuit) de
tailles différentes et un coût fixe commun.
\end{build}

\begin{omsources}
\item Commission américaine des marchés financiers (SEC), Order Competition Rule, Release 34-96495, 88 FR 128 (2023)~; Withdrawal of Proposed Regulatory Actions,
90 FR 25531 (2025).
\item Département américain de la Justice (DOJ) et Federal Trade Commission, \emph{Merger Guidelines}, 18 décembre 2023.
\item KCG Holdings, Form 8-K12G3, 1\up{er} juillet 2013~; Virtu Financial, Forms 8-K, 21 juillet 2017 et 1\up{er} mars 2019.
\item Franklin OnChain U.S. Government Money Fund, prospectus, 1\up{er} août 2026.
\item SEC, In the Matter of the Application of 24X National Exchange LLC, Release 34-101777, 27 novembre 2024~; DTCC, communiqué de presse, 29 juin 2026.
\item J. Sutton, \emph{Sunk Costs and Market Structure}, MIT Press, 1991.
\item N.~G. Mankiw et M.~D. Whinston, ``Free entry and social inefficiency'', \emph{RAND Journal of Economics}, 1986.
\end{omsources}

\section{Exercices}

\begin{exercise}[$\star$]\label{exo:fm:competition-and-the-future:1}
Calculer l'indice de six entreprises égales, puis celui des parts 33, 33, 8.5, 8.5, 8.5 et 8.5~\%.
\end{exercise}

\begin{exercise}[$\star$]\label{exo:fm:competition-and-the-future:2}
Définir un \omterm{def:fm:competition-and-the-future:moat}{avantage concurrentiel durable} et donner trois sources d'un tel avantage dans le trading.
\end{exercise}

\begin{exercise}[$\star$]\label{exo:fm:competition-and-the-future:3}
Avec $S=\$3\,000$ millions, combien d'entreprises l'entrée libre fait-elle subsister à un coût fixe de 25~millions de dollars, 50~millions de dollars et 100~millions de dollars~?
\end{exercise}

\begin{exercise}[$\star\star$]\label{exo:fm:competition-and-the-future:4}
Démontrer la \cref{prop:fm:competition-and-the-future:bounds}. Pourquoi chaque autre entreprise doit-elle être non plus grande que chacun des leaders~?
\end{exercise}

\begin{exercise}[$\star\star$]\label{exo:fm:competition-and-the-future:5}
Deux des six entreprises fusionnent sans synergie de coût marginal. La fusion serait-elle présumée réduire la concurrence selon les lignes directrices de 2023, et
est-elle rentable pour les entreprises fusionnées~?
\end{exercise}

\begin{exercise}[$\star\star$]\label{exo:fm:competition-and-the-future:6}
Comment la négociation nocturne modifie-t-elle les coûts fixes d'une entreprise de market making, et que prédit le modèle d'entrée libre pour la séance de nuit~?
\end{exercise}

\begin{exercise}[$\star\star\star$]\label{exo:fm:competition-and-the-future:7}
\emph{Programmation.} Calculer le surplus total à l'entrée libre et au nombre du planificateur pour $F=\$50$ millions. De combien l'entrée libre reste-t-elle en deçà~?
\end{exercise}

\begin{exercise}[$\star\star\star$]\label{exo:fm:competition-and-the-future:8}
\emph{Trouver la faille.} «~Les outils d'apprentissage automatique rendent la recherche bon marché, donc le métier deviendra concurrentiel et la part des leaders baissera.~»
\end{exercise}

\section{Problème~: le coût fixe baisse}

\begin{problem}[{Problème du week-end --- le coût fixe baisse}]\label{pb:fm:competition-and-the-future:1}
Le comité de stratégie d'une entreprise demande ce que feront des outils de recherche moins chers, un collatéral tokenisé et des marchés en continu aux métiers qu'elle exerce.

\medskip\noindent\textbf{Partie I --- Les concepts.}
\begin{enumerate}
\item Définir un \omterm{def:fm:competition-and-the-future:moat}{avantage concurrentiel durable}.
\item Définir la \omterm{def:fm:competition-and-the-future:hhi}{concentration du marché} et l'\omterm{def:fm:competition-and-the-future:hhi}{indice de Herfindahl--Hirschman}.
\item Définir un \omterm{def:fm:competition-and-the-future:freeentry}{équilibre à entrée libre}.
\item Qu'est-ce qui distingue les deux cas de Sutton~?
\end{enumerate}

\medskip\noindent\textbf{Partie II --- Le bilan.}
\begin{enumerate}[resume]
\item Qu'a constaté la SEC en 2022 au sujet du wholesaling, et qu'est devenue la proposition~?
\item Énoncer les seuils des lignes directrices de 2023 sur les fusions.
\item Décrire la chaîne de consolidation de la \cref{fig:fm:competition-and-the-future:chain}.
\item Qu'établissent le prospectus du \omterm{def:fm:competition-and-the-future:token}{fonds tokenisé}, l'ordre 24X et l'annonce du NSCC~?
\end{enumerate}

\medskip\noindent\textbf{Partie III --- Le modèle.}
\begin{enumerate}[resume]
\item Énoncer et démontrer la \cref{prop:fm:competition-and-the-future:bounds}~; donner les bornes pour les wholesalers.
\item Dériver le nombre d'entreprises d'entrée libre dans le modèle de Cournot symétrique.
\item Donner le nombre d'entreprises et l'indice à 25, 50 et 100~millions de dollars.
\item Donner les nombres du planificateur et expliquer l'écart.
\item Donner l'effet de la fusion sur le prix, l'indice et les profits, ainsi que la synergie qui maintient le prix.
\item Pourquoi l'indice du modèle reste-t-il en deçà des bornes publiques~?
\end{enumerate}

\medskip\noindent\textbf{Partie IV --- La stratégie.}
\begin{enumerate}[resume]
\item Lesquels des coûts fixes de l'entreprise les outils moins chers abaissent-ils, et lesquels les marchés en continu relèvent-ils~?
\item Dans lequel des métiers de l'entreprise le coût fixe est-il choisi (une course aux dépenses)~?
\item Comment savoir, à partir des chiffres propres de l'entreprise, si son avantage concurrentiel se réduit~?
\item Que faudrait-il surveiller dans le bilan au cours des deux prochaines années~?
\item Énoncer le \emph{résultat nommé}~: le nombre d'entreprises d'entrée libre et l'indice aux paramètres du livre, et la façon dont les deux évoluent lorsque le coût fixe
est divisé par deux.
\item En deux phrases, rédiger la recommandation.
\end{enumerate}
\end{problem}

\section{Questions d'entretien}

\begin{interviewq}[$\star$ \iqroles{chercheur}]\label{iq:fm:competition-and-the-future:1}
Quel est l'\omterm{def:fm:competition-and-the-future:hhi}{indice de Herfindahl--Hirschman} d'un secteur de quatre entreprises égales, et que fait à celui-ci la fusion de deux d'entre elles~?
\end{interviewq}

\begin{interviewq}[$\star$ \iqroles{trader, chercheur}]\label{iq:fm:competition-and-the-future:2}
Quel est l'avantage concurrentiel de votre entreprise, s'il en existe un~?
\end{interviewq}

\begin{interviewq}[$\star\star$ \iqroles{chercheur}]\label{iq:fm:competition-and-the-future:3}
Dans un marché de Cournot à demande linéaire et avec coût fixe, combien d'entreprises entrent-elles, et comment ce nombre varie-t-il avec la taille du marché~?
\end{interviewq}

\begin{interviewq}[$\star\star$ \iqroles{développeur}]\label{iq:fm:competition-and-the-future:4}
Qu'est-ce qui change dans un système de trading lorsque le marché échangé fonctionne 23 heures par jour~?
\end{interviewq}

\begin{interviewq}[$\star\star$ \iqroles{risque}]\label{iq:fm:competition-and-the-future:5}
Une contrepartie propose des parts d'un fonds monétaire tokenisé en collatéral un dimanche. Que vérifiez-vous~?
\end{interviewq}

\begin{interviewq}[$\star\star\star$ \iqroles{chercheur}]\label{iq:fm:competition-and-the-future:6}
Vous savez seulement que six entreprises servent un secteur et que les deux plus grandes en détiennent les deux tiers. Bornez sa concentration.
\end{interviewq}
